\chapter{Magnetic impurities, Kondo physics and screening}\label{ch:correlations}

\section{How a dot becomes a magnetic impurity}

In the local-moment regime, removing the dot's electron or adding a second one
requires much more energy than the thermal and dynamical energy scales of interest.
Charge fluctuations are suppressed, but the orientation of the remaining spin
can still fluctuate. This is the starting point for quantum magnetism in a
single impurity. An odd electron occupation supplies a spin even in a dot made
from nonmagnetic materials.

Consider first normal leads. A lead electron can tunnel onto the dot, visit the
doubly occupied configuration virtually, and tunnel back with a different spin.
A second sequence visits the empty configuration. Eliminating these high-energy
intermediate states gives an antiferromagnetic exchange Hamiltonian. For a single
effective normal conduction channel with tunnelling amplitude $t$,
\begin{equation}\label{eq:exchange-picture}
 H_K=J\bm S_d\cdot\bm s(0)+H_{\mathrm{lead}},\qquad
 J=2|t|^2\left(\frac1{-\epsilon_d}+\frac1{\epsilon_d+U}\right)>0.
\end{equation}
Here $\epsilon_d<0<\epsilon_d+U$, $\bm s(0)=
\tfrac12\sum_{\sigma\sigma'}c_\sigma\dg\bm\sigma_{\sigma\sigma'}c_{\sigma'}$
is the spin density of the coupled conduction orbital, and $\bm\sigma$ denotes
the Pauli matrices. There is also potential scattering, which vanishes at the
particle--hole-symmetric point. This low-energy reduction assumes well-separated
charge excitations. In the mixed-valence regime, retain charge fluctuations explicitly.
The normal-state effective channel is a coherent combination of the two physical
leads~\cite{hewson1993kondo,bulla2008numerical}.

Antiferromagnetic exchange favours antiparallel correlations. Its spin-flip part,
$J(S_d^+s^-+S_d^-s^+)/2$, allows the bath and impurity to exchange angular
momentum. Repeated virtual processes at progressively lower energies strengthen
the effective coupling. At the \emph{Kondo temperature} $T_K$, a weak-exchange
description fails and the low-energy state becomes strongly correlated. For
a spin-$1/2$ impurity with one screening channel and a metallic bath, the
zero-temperature ground state is a singlet. This is Kondo screening.

\section{Physical mechanisms of screening}\label{sec:screening-picture}

The word screening is used for several different physical processes. Ordinary
electrostatic screening rearranges charge around a perturbing electric potential.
Kondo screening involves quantum \emph{spin} correlations. Proximity pairing
mixes different even-charge configurations through exchange of Cooper pairs.
\Cref{fig:primer-screening} distinguishes these mechanisms.

\begin{figure}[htb]
\centering\manualfigure{primer_screening}
\caption{Three distinct physical mechanisms. Electrostatic screening
rearranges charge; proximity pairing coherently mixes the empty and doubly occupied
dot; Kondo screening entangles a local moment with conduction electrons. The
screening cloud schematically represents correlations with many conduction electrons.}
\label{fig:primer-screening}
\end{figure}

The two-spin singlet
\begin{equation}
 \ket{\mathrm{singlet}}=\frac{\ket{\up_d\dn_b}-\ket{\dn_d\up_b}}{\sqrt2}
\end{equation}
is a useful illustration of entanglement. Its total spin is zero, each local
spin has zero mean polarization, but each has $\expval{\bm S_d^2}=3/4$ and
$\expval{\bm S_d\cdot\bm S_b}=-3/4$: the local moments remain finite and are
antiferromagnetically correlated. This two-spin wavefunction illustrates the entanglement.
In the full Kondo problem the bath partner is a many-electron state distributed
across energies and positions. A normal-metal correlation length is of order
$\xi_K\sim\hbar v_F/(k_BT_K)$, with its precise definition depending on the
correlation being measured.

A singlet in a proximitized dot can instead take the form $u\ket0+v\ket2$. In the strict
large-gap model it contains exclusively empty and doubly occupied configurations. Both states have
total spin zero, but only the first picture screens a pre-existing local moment.
At finite gap and finite hybridization these descriptions can evolve smoothly
into one another because they share the same symmetries, allowing a crossover
between pairing-dominated and Kondo-dominated singlet regimes.

\begin{physics}[Which observable establishes screening?]
To identify the screening mechanism, examine the dot charge distribution,
$\expval{\bm S_d^2}$, spin correlations with the bath,
the temperature dependence of the susceptibility, and the normal-state Kondo
scale. These distinguish a local moment, its screening and a predominantly
empty/doubly occupied proximity singlet. Zero local magnetization,
$\expval{S_d^z}=0$, is compatible with all of these and with an unpolarized
ensemble of free spins; a total-spin label $S=0$ likewise leaves the mechanism to be identified.
\end{physics}

\section{Why a superconducting gap competes with Kondo screening}

The Kondo effect draws on low-energy electron--hole excitations of a normal
metal. A superconducting gap removes the gapless continuum on which this
low-energy buildup relies. In the local-moment regime the useful comparison is
therefore between $k_BT_K$, defined for the corresponding \emph{normal-state}
impurity, and $\Delta$~\cite{choi2004kondo,siano2004josephson,bauer2007spectral}.

When $k_BT_K\gg\Delta$, correlations have become strong before the gap cuts off
the normal-metal flow. A screened singlet is favoured. When $k_BT_K\ll\Delta$,
the gap interrupts that flow before screening is complete, and a doublet is
favoured. When the two scales are comparable, locating the transition accurately
requires an explicit solution of the superconducting impurity problem.

\begin{figure}[htb]
\centering\manualfigure{primer_scales}
\caption{Schematic energy-scale comparison in the local-moment regime. Following
the normal-state flow downwards, strong Kondo correlations are reached before the
gap in the left panel, whereas the gap intervenes first in the right panel.
Positions on these logarithmic energy scales illustrate the two limiting regimes.
The shaded regions denote energies unavailable as bulk single-QP
excitations.}
\label{fig:primer-scales}
\end{figure}

The Kondo scale is exponentially sensitive to microscopic parameters. Definitions
from susceptibility, a spectral width and transport fits differ by numerical
factors. Total and per-lead definitions of $\Gamma$ introduce another common
conversion. Consequently, a quoted critical value of $T_K/\Delta$ is meaningful
only with its scale convention, phase, asymmetry and Hamiltonian specified. This
manual writes $k_BT_K/\Delta$ when the Kondo scale is expressed as a temperature; setting $k_B=1$
converts it to an energy. Determine the transition boundary for the specified
junction and conventions.

Phase matters even at fixed normal-state $T_K$. Contributions from the two
condensates interfere, and their cancellation near $\phi=\pi$ can destabilize a
singlet. Detuning can remove the local moment entirely. Thus a diagram with axes
$T_K$ and $\Delta$ summarizes one competition, while gate--phase diagrams often
describe experiments more directly. Emptying the dot produces a singlet through
a change of charge configuration, with the local moment removed.

\section{The classical Yu--Shiba--Rusinov limit}\label{sec:classical-ysr}

A different instructive limit replaces the impurity by a prescribed classical
spin $\bm S$, with no quantum spin flips. Magnetic scattering affects the two
members of a Cooper pair differently and binds a localized state. In a
particle--hole-symmetric, wide-band $s$-wave host, with no potential scattering,
one signed YSR pole can be written~\cite{yu1965bound,shiba1968classical,rusinov1969superconductivity,balatsky2006impurity}
\begin{equation}\label{eq:classical-ysr}
 E_{\mathrm{YSR}}=\Delta\frac{1-\alpha^2}{1+\alpha^2},\qquad
 \alpha=\frac{\pi\rho_0JS}{2}.
\end{equation}
Here $\rho_0$ is the normal density of states per spin and the exchange is
$J\bm S\cdot c\dg\bm\sigma c/2$; these conventions fix the factor of two.
The particle--hole partner has the opposite pole energy. Potential scattering
changes the formula and the electron/hole weights.

\begin{figure}[htb]
\centering\manualfigure{primer_ysr}
\caption{Classical-spin YSR poles from \cref{eq:classical-ysr}. (a) Following a
signed branch through zero and its particle--hole partner. (b) The positive
excitation cost $|E_{\mathrm{YSR}}|$. The change of sign changes the reference
occupation while the excitation cost remains nonnegative. These curves describe
scattering from a prescribed classical spin.}
\label{fig:primer-ysr}
\end{figure}

Starting at weak exchange, the pole leaves the gap edge and moves towards zero.
At $\alpha=1$ its sign changes, requiring a change in the occupation defining the
many-body ground state. Excitations above that ground state cost $|E_{\mathrm{YSR}}|$.
The binding energy below the continuum edge is the separate quantity
$\Delta-|E_{\mathrm{YSR}}|$.

The classical model explains localized impurity states and a parity-changing
crossing within a simple analytical description based on static magnetic scattering.
Describing the Kondo effect requires quantum impurity-spin dynamics. When a
quantum-dot spectrum resembles the classical curve, examine its many-body
wavefunctions and spin correlations to identify the screening mechanism.

\section{Quantum YSR states and the singlet--doublet transition}

For a quantum impurity, diagonalization produces many-body states of the
\emph{dot plus superconducting leads}. In the spin-rotation-invariant,
single-orbital problem, the relevant low states are usually an even singlet $S$
and an odd doublet $D$. An electron probe connects them. Define the signed
energy difference
\begin{equation}\label{eq:sd-energy}
 \delta_{SD}=E_D-E_S.
\end{equation}
The singlet is the ground state for $\delta_{SD}>0$, and the doublet for
$\delta_{SD}<0$. The lowest transition between them costs $|\delta_{SD}|$;
positive and negative bias probe electron addition and removal, with generally
different weights. A branch crossing through zero in a plot of $\delta_{SD}$
and two resonances meeting at zero bias are related representations of the
same singlet--doublet transition.

Parity conservation keeps the even and odd sectors decoupled, allowing their
levels to cross exactly. Coupled states of the same symmetry instead form an
avoided crossing. At zero temperature the parity crossing changes the global
ground state. A closed, isolated system retains its initially prepared parity
as a parameter is varied through the crossing. Relaxation to the new global ground state requires a channel for
parity exchange with the surroundings.

\begin{figure}[htb]
\centering\manualfigure{primer_thermal}
\caption{Ground-state competition and thermal occupation of a singlet and a
spin-degenerate doublet. The horizontal coordinate is the actual energy difference
$\delta_{SD}$ in units of $k_BT$. The doublet
probability assumes equilibrium between both parities and no other thermally
occupied states. Thermal rounding changes the populations while the
Hamiltonian eigenvalues retain their exact crossing.}
\label{fig:primer-thermal}
\end{figure}

At temperatures low enough that only these multiplets are appreciably populated, equilibrium gives
\begin{equation}\label{eq:doublet-probability}
 Z=\ee^{-\beta E_S}+2\ee^{-\beta E_D},\qquad
 P_D=\frac{2\ee^{-\beta\delta_{SD}}}{1+2\ee^{-\beta\delta_{SD}}}.
\end{equation}
The factor two is the doublet degeneracy. At equal energies the doublet has
probability $2/3$. A magnetic field splits the two Boltzmann weights;
a parity-trapped experiment instead needs a restricted population model. The
parity crossing remains exact as temperature changes the populations and
spectral broadening changes the displayed line shapes.

In this setting, \emph{Andreev bound state} is often used broadly for a subgap
excitation of a superconducting weak link, while \emph{YSR state} emphasizes
its magnetic-impurity origin. For an interacting quantum dot these names often
refer to the same measured spectral features. We specify the state, transition
and approximation to make the physical interpretation explicit. A useful
short introduction to this many-body interpretation is
Ref.~\cite{zitko2018superconducting}; the wider theory is reviewed in
Refs.~\cite{martinrodero2011josephson,meden2019andersonjosephson}.

\section{A residual doublet can have a reduced local spin}

A doublet has total spin $1/2$, distributed between dot and leads.
For its $S_z=+1/2$ component,
\begin{equation}
 \expval{S_d^z}+\expval{S_{\mathrm{leads}}^z}=\frac12.
\end{equation}
Hybridization changes the local charge fluctuations and dot--lead correlations,
redistributing this polarization. Calling a doublet ``partially screened''
describes this local reduction while the conserved total spin remains $1/2$.

If a weak field acts only on the dot as $B S_d^z$, the Hellmann--Feynman theorem
relates the linear splitting to the local polarization. In an SU(2)-symmetric
zero-field doublet,
\begin{equation}\label{eq:local-spin-slope}
 \lim_{B\to0}\frac{|E_{D_\up}-E_{D_\dn}|}{|B|}
 =2\bigl|\matel{D_\up}{S_d^z}{D_\up}\bigr|.
\end{equation}
When $B=g\mu_B B_{\mathrm{lab}}$, this slope is $|g_{\mathrm{eff}}/g|$.
It is an impurity Knight-shift diagnostic~\cite{pavesic2023impurity}. A field
coupled uniformly to the \emph{total} spin would instead measure the total-spin
response. Spin--orbit coupling and different lead $g$-factors require a more
general analysis, including directional dependence.

\Cref{ch:ex-esr} demonstrates the same local-spin renormalization in the Zeeman
splitting and in the matrix element for a transverse magnetic drive. A reduced
$g_{\mathrm{eff}}$ is experimentally relevant to a spin qubit. Interpreting it
quantitatively requires accounting for charge admixture and wavefunction
delocalization alongside Kondo correlations, since each can reduce the local polarization.

\begin{physics}[Read the figures]
A gate sweep takes a tunnel resonance down to zero energy and back up. Which
physical change could produce this trace? It can cross a singlet--doublet boundary
as charge or induced pairing changes. Identify the normal-state regime, the field
response and the transition weights to assess the role of Kondo correlations. Likewise,
observing a doublet with a reduced dot-local Zeeman slope is compatible with
total spin $1/2$; it reveals where that spin responds to the applied field.
\end{physics}

\paragraph{Reading further.}
Hewson's book~\cite{hewson1993kondo} develops the normal-state local-moment and
Kondo problems. Bulla, Costi and Pruschke~\cite{bulla2008numerical} explain why
many energy scales must be retained and how NRG does so. Balatsky, Vekhter and
Zhu~\cite{balatsky2006impurity} review impurity-induced bound states, while
Heinrich, Pascual and Franke~\cite{heinrich2018single} connect individual
magnetic impurities to superconducting-surface experiments.
